\documentclass[12pt]{article}
\usepackage[T1]{fontenc}
\usepackage[utf8]{inputenc}
\usepackage{lmodern}
\usepackage{textcomp}
\usepackage{enumitem}
\usepackage{geometry}
\usepackage{graphicx}
\usepackage{amsmath, amssymb}
\geometry{margin=1in}
\begin{document}
\begin{center}
Chemistry - Undergraduate Level Assessment 17 \
Total Marks: 40 \
Answer all questions.
\end{center}

\section*{Multiple Choice (10 marks)}
\begin{enumerate}[label=\arabic*.]
\item Which of the following is an n-type semiconductor?
\begin{enumerate}[label=(\alph*)]
    \item Silicon
    \item Diamond
    \item Silicon carbide
    \item Arsenic-doped silicon
\end{enumerate}
\item Calculate the polarization of a proton in a magnetic field of 335 mT and 10.5 T at 298 K.
\begin{enumerate}[label=(\alph*)]
    \item 6.345 x 10^-4 at 0.335 T; 9.871 x 10^-5 at 10.5 T
    \item 0.793 x 10^-4 at 0.335 T; 6.931 x 10^-7 at 10.5 T
    \item 1.148 x 10^-6 at 0.335 T; 3.598 x 10^-5 at 10.5 T
    \item 4.126 x 10^-3 at 0.335 T; 2.142 x 10^-6 at 10.5 T
\end{enumerate}
\item Which of the following is lower for argon than for neon?
\begin{enumerate}[label=(\alph*)]
    \item Melting point
    \item Boiling point
    \item Polarizability
    \item First ionization energy
\end{enumerate}
\item Of the following atoms, which has the lowest electron affinity?
\begin{enumerate}[label=(\alph*)]
    \item F
    \item Si
    \item O
    \item Ca
\end{enumerate}
\item Of the following compounds, which has the lowest melting point?
\begin{enumerate}[label=(\alph*)]
    \item HCl
    \item AgCl
    \item CaCl 2
    \item CCl 4
\end{enumerate}
\end{enumerate}

\section*{True / False (10 marks)}
\textit{Answer True or False}
\begin{enumerate}[label=\arabic*.]
\setcounter{enumi}{5}
\item True or False: The balanced equation shows a MnO 4- to I- ratio of 6:5.
\item True or False: Nuclear magnetic resonance spectroscopy is a light-scattering technique that analyzes molecular structures.
\item True or False: CH 3 F has a 1 H resonance that is closer to TMS than CH 3 Cl on a spectrometer with a 12.0 T magnet.
\item True or False: Protons in a molecule with a rotational correlation time of 1 ns have the fastest spin-lattice relaxation at a magnetic field of 3.74 T.
\item True or False: The 55 Mn nuclide has five allowed energy levels for its nuclear states.
\end{enumerate}

\section*{Long Form Response (20 marks)}
\textit{Answer each question with a clear and complete explanation.}
\begin{enumerate}[label=\arabic*.]
\setcounter{enumi}{10}
\item Discuss the factors influencing the ratio of line intensities in the EPR spectrum of the t-Bu radical (CH 3)3 C• and explain the significance of these ratios in understanding radical behavior.
\item Discuss the significance of NMR frequency in relation to nuclear spin properties, and explain why 17 O exhibits a frequency of 115.5 MHz in a 20.0 T magnetic field.
\end{enumerate}

\vfill
\noindent\textit{End of Paper}
\end{document}